% ---------------------------------------------------------------------------
% Implementation of Algorithms 1, 2 and 3 of Duan et al. (2025).
% Content of the existing \subsection{Implementation details}: put
% % ---------------------------------------------------------------------------
% Implementation of Algorithms 1, 2 and 3 of Duan et al. (2025).
% Content of the existing \subsection{Implementation details}: put
% % ---------------------------------------------------------------------------
% Implementation of Algorithms 1, 2 and 3 of Duan et al. (2025).
% Content of the existing \subsection{Implementation details}: put
% % ---------------------------------------------------------------------------
% Implementation of Algorithms 1, 2 and 3 of Duan et al. (2025).
% Content of the existing \subsection{Implementation details}: put
% \input{implementation_details} right below that heading.
% Preamble packages required:
%   \usepackage{amsmath}
%   \usepackage{listings}
%   \usepackage{xcolor}
% Citation key used: duan2025 (entry in implementation_details.bib).
% ---------------------------------------------------------------------------

\lstdefinestyle{pycode}{
  language=Python,
  basicstyle=\ttfamily\footnotesize,
  keywordstyle=\bfseries\color{blue!60!black},
  commentstyle=\itshape\color{green!40!black},
  stringstyle=\color{red!60!black},
  numbers=left, numberstyle=\tiny\color{gray}, numbersep=6pt,
  frame=single, breaklines=true, showstringspaces=false,
  columns=fullflexible, keepspaces=true, tabsize=4
}

The three algorithms of Duan et al.~\cite{duan2025} are implemented in Python as methods of
the class \texttt{DuanSSSP}, in the file \texttt{duan.py}: \texttt{find\_pivots}
(Algorithm~1), \texttt{base\_case} (Algorithm~2) and \texttt{bmssp} (Algorithm~3). The
listings below are the code of these methods. Only the lines that write logs and update the
operation counters of the experiment were removed, since they do not change the result. We
checked that this version returns the same distances and predecessors as the full one.

\paragraph{Shared state.} The methods work on four lists indexed by vertex:
\texttt{d[v]} holds the estimate $\hat{d}[v]$, \texttt{h[v]} the number of arcs of the
current path to $v$, \texttt{pred[v]} the predecessor $\mathrm{Pred}[v]$, and
\texttt{pred\_w[v]} the weight of the arc $(\mathrm{Pred}[v], v)$. Adjacency lists are flat
tuples $(v_0, w_0, v_1, w_1, \dots)$, read in pairs with \texttt{zip(adj\_u, adj\_u)}. The
parameter \texttt{depth} is only used to indent the logs.

\paragraph{Total order (Assumption 2.1).} The paper assumes that all paths have distinct
lengths. We represent the value of $\hat{d}[v]$ by the tuple $(\hat{d}[v], h[v], v)$, which
Python compares lexicographically, so every bound $B$, $B_i$ and $B'_i$ is a tuple of this
form ($B=\infty$ is $(\infty,\infty,\infty)$). The relaxation test of Remark~3.4,
$\hat{d}[u]+w_{uv}\le\hat{d}[v]$, becomes
\begin{equation*}
  \bigl(\hat{d}[u]+w,\; h[u]+1,\; u\bigr) \;\le\; \bigl(\hat{d}[v],\; h[v],\; \mathrm{Pred}[v]\bigr),
\end{equation*}
because the two paths compared end at $v$ and the first possible difference after the
length and the number of arcs is the predecessor. The ``$\le$'' lets an arc that is already
tight be relaxed again, as the paper requires. This test is written directly inside each of
the three methods.

\subsubsection{Algorithm 1: FindPivots}

\texttt{find\_pivots(B, S)} (Listing~\ref{lst:findpivots}) decides which vertices of $S$
need to be passed on to the recursion.
\begin{enumerate}
  \item It runs $k$ rounds of Bellman--Ford starting from $S$. In each round only the
        vertices reached in the previous round (\texttt{prev}) are expanded, and a vertex
        enters $W_i$ only if its new value is below $B$.
  \item If $|W|$ exceeds $k|S|$, it stops and returns $P=S$: the frontier is too large to
        be reduced.
  \item Otherwise it builds the forest $F$ of tight arcs inside $W$. Under the total order
        a vertex has at most one tight incoming arc, the one from $\mathrm{Pred}[v]$, so
        $F$ is read from \texttt{pred} in $O(|W|)$ time. The two equalities check that the
        arc is still tight, because $\mathrm{Pred}[v]$ may be outdated if $u$ was improved
        later.
  \item The pivots are the roots in $S$ whose tree in $F$ has at least $k$ vertices. Tree
        sizes are computed with an explicit-stack DFS.
\end{enumerate}

\begin{lstlisting}[style=pycode, caption={Algorithm 1, FindPivots (\texttt{duan.py}).}, label={lst:findpivots}]
def find_pivots(self, B, S, depth):
    k, adj = self.k, self.g.adj
    d, h, pred, pred_w = self.d, self.h, self.pred, self.pred_w
    W = set(S)
    prev = set(S)
    for i in range(1, k + 1):
        Wi = set()
        for u in prev:
            adj_u = iter(adj[u])
            for v, w in zip(adj_u, adj_u):
                du, hu = d[u] + w, h[u] + 1
                if (du, hu, u) <= (d[v], h[v], pred[v]):
                    d[v], h[v], pred[v], pred_w[v] = du, hu, u, w
                    if (du, hu, v) < B:
                        Wi.add(v)
        W |= Wi
        if len(W) > k * len(S):
            return set(S), W
        prev = Wi

    children = {u: [] for u in W}
    has_parent = set()
    for v in W:
        u = pred[v]
        if u in children and d[v] == d[u] + pred_w[v] and h[v] == h[u] + 1:
            children[u].append(v)
            has_parent.add(v)

    P = set()
    for r in S:
        if r in has_parent:
            continue
        size, stack = 0, [r]
        while stack:
            x = stack.pop()
            size += 1
            stack.extend(children[x])
        if size >= k:
            P.add(r)
    return P, W
\end{lstlisting}

\subsubsection{Algorithm 2: BaseCase}

\texttt{base\_case(B, S)} (Listing~\ref{lst:basecase}) is the bottom of the recursion
($l=0$), where $S=\{x\}$ is a single vertex.
\begin{enumerate}
  \item It runs a Dijkstra from $x$ that settles at most $k+1$ vertices and relaxes only
        paths whose value is below $B$.
  \item The paper uses a heap with DecreaseKey. We use Python's binary heap
        (\texttt{heapq}) with lazy deletion: an improved vertex is pushed again, and
        outdated entries are skipped when popped, by comparing the stored key with the
        current one. \texttt{heap\_ops} returns the push and pop functions of the heap
        selected for the run.
  \item If at most $k$ vertices were settled, the execution is successful and returns
        $B'=B$ and all of them. Otherwise the largest value among the $k+1$ settled
        vertices becomes the new bound $B'$, and only the vertices below it are returned.
\end{enumerate}

\begin{lstlisting}[style=pycode, caption={Algorithm 2, BaseCase (\texttt{duan.py}).}, label={lst:basecase}]
def base_case(self, B, S, depth):
    d, h, pred, pred_w = self.d, self.h, self.pred, self.pred_w
    adj = self.g.adj
    (x,) = S
    U0 = {x}
    push, pop = heap_ops(self.contar)
    heap = []
    kx = (d[x], h[x], x)
    push(heap, (kx, x))
    in_heap = {x: kx}
    done = set()
    while heap and len(U0) < self.k + 1:
        kv, u = pop(heap)
        if u in done or kv != (d[u], h[u], u):
            continue
        done.add(u)
        U0.add(u)
        adj_u = iter(adj[u])
        for v, w in zip(adj_u, adj_u):
            du, hu = d[u] + w, h[u] + 1
            if (du, hu, u) <= (d[v], h[v], pred[v]) and (du, hu, v) < B:
                d[v], h[v], pred[v], pred_w[v] = du, hu, u, w
                kv2 = (du, hu, v)
                old = in_heap.get(v)
                if old is None or kv2 < old:
                    in_heap[v] = kv2
                    push(heap, (kv2, v))
    if len(U0) <= self.k:
        return B, U0
    Bp = max((d[v], h[v], v) for v in U0)
    U = {v for v in U0 if (d[v], h[v], v) < Bp}
    return Bp, U
\end{lstlisting}

\subsubsection{Algorithm 3: BMSSP}

\texttt{bmssp(l, B, S)} (Listing~\ref{lst:bmssp}) is the recursive core. The solver starts
with \texttt{bmssp(l, (inf, inf, inf), \{s\})} at the top level $l$.
\begin{enumerate}
  \item At $l=0$ it delegates to \texttt{base\_case}. Otherwise it calls
        \texttt{find\_pivots} and inserts the pivots into the block-based structure $D$ of
        Lemma~3.3 (class \texttt{BlockDS}), with block size $M=2^{(l-1)t}$. $B'_0$ is the
        smallest value among the pivots.
  \item While $|U| < k\cdot 2^{lt}$ and $D$ is not empty, it pulls from $D$ the (up to)
        $M$ smallest keys $S_i$ with a separating bound $B_i$, and recurses with
        \texttt{bmssp(l-1, B\_i, S\_i)}, which returns $(B'_i, U_i)$.
  \item It relaxes the arcs leaving $U_i$. A relaxed vertex whose new value lies in
        $[B_i, B)$ goes back into $D$ with Insert. One in $[B'_i, B_i)$ goes into the list
        \texttt{K}, together with the vertices of $S_i$ that the recursive call did not
        complete. \texttt{K} is returned to $D$ with a single BatchPrepend, since all its
        values are smaller than every value still in $D$.
  \item At the end, $B'=\min(B'_i, B)$, and $U$ also receives the vertices of $W$ (from
        FindPivots) whose value is below $B'$. $B'=B$ means a successful execution;
        $B'<B$ means it stopped early because of a large workload.
\end{enumerate}
The paper checks the workload condition inside the loop; here it is part of the
\texttt{while} condition, which has the same effect.

\begin{lstlisting}[style=pycode, caption={Algorithm 3, BMSSP (\texttt{duan.py}).}, label={lst:bmssp}]
def bmssp(self, l, B, S, depth=0):
    d, h, pred, pred_w = self.d, self.h, self.pred, self.pred_w
    adj = self.g.adj
    if l == 0:
        return self.base_case(B, S, depth)

    P, W = self.find_pivots(B, S, depth)
    M = 2 ** ((l - 1) * self.t)
    D = BlockDS(M, B, self.fmt, depth, self.contar)
    for x in P:
        D.insert(x, (d[x], h[x], x))
    Bp_last = min(((d[x], h[x], x) for x in P), default=B)
    U = set()
    limit = self.k * 2 ** (l * self.t)

    while len(U) < limit and len(D) > 0:
        Bi, Si = D.pull()
        Bpi, Ui = self.bmssp(l - 1, Bi, set(Si), depth + 1)
        U |= Ui

        K = []
        for u in Ui:
            adj_u = iter(adj[u])
            for v, w in zip(adj_u, adj_u):
                du, hu = d[u] + w, h[u] + 1
                if (du, hu, u) <= (d[v], h[v], pred[v]):
                    d[v], h[v], pred[v], pred_w[v] = du, hu, u, w
                    val = (du, hu, v)
                    if Bi <= val < B:
                        D.insert(v, val)
                    elif Bpi <= val < Bi:
                        K.append((v, val))
        K.extend((x, (d[x], h[x], x)) for x in Si if Bpi <= (d[x], h[x], x) < Bi)
        D.batch_prepend(K)
        Bp_last = Bpi

    Bp = min(Bp_last, B)
    U |= {x for x in W if (d[x], h[x], x) < Bp}
    return Bp, U
\end{lstlisting}
 right below that heading.
% Preamble packages required:
%   \usepackage{amsmath}
%   \usepackage{listings}
%   \usepackage{xcolor}
% Citation key used: duan2025 (entry in implementation_details.bib).
% ---------------------------------------------------------------------------

\lstdefinestyle{pycode}{
  language=Python,
  basicstyle=\ttfamily\footnotesize,
  keywordstyle=\bfseries\color{blue!60!black},
  commentstyle=\itshape\color{green!40!black},
  stringstyle=\color{red!60!black},
  numbers=left, numberstyle=\tiny\color{gray}, numbersep=6pt,
  frame=single, breaklines=true, showstringspaces=false,
  columns=fullflexible, keepspaces=true, tabsize=4
}

The three algorithms of Duan et al.~\cite{duan2025} are implemented in Python as methods of
the class \texttt{DuanSSSP}, in the file \texttt{duan.py}: \texttt{find\_pivots}
(Algorithm~1), \texttt{base\_case} (Algorithm~2) and \texttt{bmssp} (Algorithm~3). The
listings below are the code of these methods. Only the lines that write logs and update the
operation counters of the experiment were removed, since they do not change the result. We
checked that this version returns the same distances and predecessors as the full one.

\paragraph{Shared state.} The methods work on four lists indexed by vertex:
\texttt{d[v]} holds the estimate $\hat{d}[v]$, \texttt{h[v]} the number of arcs of the
current path to $v$, \texttt{pred[v]} the predecessor $\mathrm{Pred}[v]$, and
\texttt{pred\_w[v]} the weight of the arc $(\mathrm{Pred}[v], v)$. Adjacency lists are flat
tuples $(v_0, w_0, v_1, w_1, \dots)$, read in pairs with \texttt{zip(adj\_u, adj\_u)}. The
parameter \texttt{depth} is only used to indent the logs.

\paragraph{Total order (Assumption 2.1).} The paper assumes that all paths have distinct
lengths. We represent the value of $\hat{d}[v]$ by the tuple $(\hat{d}[v], h[v], v)$, which
Python compares lexicographically, so every bound $B$, $B_i$ and $B'_i$ is a tuple of this
form ($B=\infty$ is $(\infty,\infty,\infty)$). The relaxation test of Remark~3.4,
$\hat{d}[u]+w_{uv}\le\hat{d}[v]$, becomes
\begin{equation*}
  \bigl(\hat{d}[u]+w,\; h[u]+1,\; u\bigr) \;\le\; \bigl(\hat{d}[v],\; h[v],\; \mathrm{Pred}[v]\bigr),
\end{equation*}
because the two paths compared end at $v$ and the first possible difference after the
length and the number of arcs is the predecessor. The ``$\le$'' lets an arc that is already
tight be relaxed again, as the paper requires. This test is written directly inside each of
the three methods.

\subsubsection{Algorithm 1: FindPivots}

\texttt{find\_pivots(B, S)} (Listing~\ref{lst:findpivots}) decides which vertices of $S$
need to be passed on to the recursion.
\begin{enumerate}
  \item It runs $k$ rounds of Bellman--Ford starting from $S$. In each round only the
        vertices reached in the previous round (\texttt{prev}) are expanded, and a vertex
        enters $W_i$ only if its new value is below $B$.
  \item If $|W|$ exceeds $k|S|$, it stops and returns $P=S$: the frontier is too large to
        be reduced.
  \item Otherwise it builds the forest $F$ of tight arcs inside $W$. Under the total order
        a vertex has at most one tight incoming arc, the one from $\mathrm{Pred}[v]$, so
        $F$ is read from \texttt{pred} in $O(|W|)$ time. The two equalities check that the
        arc is still tight, because $\mathrm{Pred}[v]$ may be outdated if $u$ was improved
        later.
  \item The pivots are the roots in $S$ whose tree in $F$ has at least $k$ vertices. Tree
        sizes are computed with an explicit-stack DFS.
\end{enumerate}

\begin{lstlisting}[style=pycode, caption={Algorithm 1, FindPivots (\texttt{duan.py}).}, label={lst:findpivots}]
def find_pivots(self, B, S, depth):
    k, adj = self.k, self.g.adj
    d, h, pred, pred_w = self.d, self.h, self.pred, self.pred_w
    W = set(S)
    prev = set(S)
    for i in range(1, k + 1):
        Wi = set()
        for u in prev:
            adj_u = iter(adj[u])
            for v, w in zip(adj_u, adj_u):
                du, hu = d[u] + w, h[u] + 1
                if (du, hu, u) <= (d[v], h[v], pred[v]):
                    d[v], h[v], pred[v], pred_w[v] = du, hu, u, w
                    if (du, hu, v) < B:
                        Wi.add(v)
        W |= Wi
        if len(W) > k * len(S):
            return set(S), W
        prev = Wi

    children = {u: [] for u in W}
    has_parent = set()
    for v in W:
        u = pred[v]
        if u in children and d[v] == d[u] + pred_w[v] and h[v] == h[u] + 1:
            children[u].append(v)
            has_parent.add(v)

    P = set()
    for r in S:
        if r in has_parent:
            continue
        size, stack = 0, [r]
        while stack:
            x = stack.pop()
            size += 1
            stack.extend(children[x])
        if size >= k:
            P.add(r)
    return P, W
\end{lstlisting}

\subsubsection{Algorithm 2: BaseCase}

\texttt{base\_case(B, S)} (Listing~\ref{lst:basecase}) is the bottom of the recursion
($l=0$), where $S=\{x\}$ is a single vertex.
\begin{enumerate}
  \item It runs a Dijkstra from $x$ that settles at most $k+1$ vertices and relaxes only
        paths whose value is below $B$.
  \item The paper uses a heap with DecreaseKey. We use Python's binary heap
        (\texttt{heapq}) with lazy deletion: an improved vertex is pushed again, and
        outdated entries are skipped when popped, by comparing the stored key with the
        current one. \texttt{heap\_ops} returns the push and pop functions of the heap
        selected for the run.
  \item If at most $k$ vertices were settled, the execution is successful and returns
        $B'=B$ and all of them. Otherwise the largest value among the $k+1$ settled
        vertices becomes the new bound $B'$, and only the vertices below it are returned.
\end{enumerate}

\begin{lstlisting}[style=pycode, caption={Algorithm 2, BaseCase (\texttt{duan.py}).}, label={lst:basecase}]
def base_case(self, B, S, depth):
    d, h, pred, pred_w = self.d, self.h, self.pred, self.pred_w
    adj = self.g.adj
    (x,) = S
    U0 = {x}
    push, pop = heap_ops(self.contar)
    heap = []
    kx = (d[x], h[x], x)
    push(heap, (kx, x))
    in_heap = {x: kx}
    done = set()
    while heap and len(U0) < self.k + 1:
        kv, u = pop(heap)
        if u in done or kv != (d[u], h[u], u):
            continue
        done.add(u)
        U0.add(u)
        adj_u = iter(adj[u])
        for v, w in zip(adj_u, adj_u):
            du, hu = d[u] + w, h[u] + 1
            if (du, hu, u) <= (d[v], h[v], pred[v]) and (du, hu, v) < B:
                d[v], h[v], pred[v], pred_w[v] = du, hu, u, w
                kv2 = (du, hu, v)
                old = in_heap.get(v)
                if old is None or kv2 < old:
                    in_heap[v] = kv2
                    push(heap, (kv2, v))
    if len(U0) <= self.k:
        return B, U0
    Bp = max((d[v], h[v], v) for v in U0)
    U = {v for v in U0 if (d[v], h[v], v) < Bp}
    return Bp, U
\end{lstlisting}

\subsubsection{Algorithm 3: BMSSP}

\texttt{bmssp(l, B, S)} (Listing~\ref{lst:bmssp}) is the recursive core. The solver starts
with \texttt{bmssp(l, (inf, inf, inf), \{s\})} at the top level $l$.
\begin{enumerate}
  \item At $l=0$ it delegates to \texttt{base\_case}. Otherwise it calls
        \texttt{find\_pivots} and inserts the pivots into the block-based structure $D$ of
        Lemma~3.3 (class \texttt{BlockDS}), with block size $M=2^{(l-1)t}$. $B'_0$ is the
        smallest value among the pivots.
  \item While $|U| < k\cdot 2^{lt}$ and $D$ is not empty, it pulls from $D$ the (up to)
        $M$ smallest keys $S_i$ with a separating bound $B_i$, and recurses with
        \texttt{bmssp(l-1, B\_i, S\_i)}, which returns $(B'_i, U_i)$.
  \item It relaxes the arcs leaving $U_i$. A relaxed vertex whose new value lies in
        $[B_i, B)$ goes back into $D$ with Insert. One in $[B'_i, B_i)$ goes into the list
        \texttt{K}, together with the vertices of $S_i$ that the recursive call did not
        complete. \texttt{K} is returned to $D$ with a single BatchPrepend, since all its
        values are smaller than every value still in $D$.
  \item At the end, $B'=\min(B'_i, B)$, and $U$ also receives the vertices of $W$ (from
        FindPivots) whose value is below $B'$. $B'=B$ means a successful execution;
        $B'<B$ means it stopped early because of a large workload.
\end{enumerate}
The paper checks the workload condition inside the loop; here it is part of the
\texttt{while} condition, which has the same effect.

\begin{lstlisting}[style=pycode, caption={Algorithm 3, BMSSP (\texttt{duan.py}).}, label={lst:bmssp}]
def bmssp(self, l, B, S, depth=0):
    d, h, pred, pred_w = self.d, self.h, self.pred, self.pred_w
    adj = self.g.adj
    if l == 0:
        return self.base_case(B, S, depth)

    P, W = self.find_pivots(B, S, depth)
    M = 2 ** ((l - 1) * self.t)
    D = BlockDS(M, B, self.fmt, depth, self.contar)
    for x in P:
        D.insert(x, (d[x], h[x], x))
    Bp_last = min(((d[x], h[x], x) for x in P), default=B)
    U = set()
    limit = self.k * 2 ** (l * self.t)

    while len(U) < limit and len(D) > 0:
        Bi, Si = D.pull()
        Bpi, Ui = self.bmssp(l - 1, Bi, set(Si), depth + 1)
        U |= Ui

        K = []
        for u in Ui:
            adj_u = iter(adj[u])
            for v, w in zip(adj_u, adj_u):
                du, hu = d[u] + w, h[u] + 1
                if (du, hu, u) <= (d[v], h[v], pred[v]):
                    d[v], h[v], pred[v], pred_w[v] = du, hu, u, w
                    val = (du, hu, v)
                    if Bi <= val < B:
                        D.insert(v, val)
                    elif Bpi <= val < Bi:
                        K.append((v, val))
        K.extend((x, (d[x], h[x], x)) for x in Si if Bpi <= (d[x], h[x], x) < Bi)
        D.batch_prepend(K)
        Bp_last = Bpi

    Bp = min(Bp_last, B)
    U |= {x for x in W if (d[x], h[x], x) < Bp}
    return Bp, U
\end{lstlisting}
 right below that heading.
% Preamble packages required:
%   \usepackage{amsmath}
%   \usepackage{listings}
%   \usepackage{xcolor}
% Citation key used: duan2025 (entry in implementation_details.bib).
% ---------------------------------------------------------------------------

\lstdefinestyle{pycode}{
  language=Python,
  basicstyle=\ttfamily\footnotesize,
  keywordstyle=\bfseries\color{blue!60!black},
  commentstyle=\itshape\color{green!40!black},
  stringstyle=\color{red!60!black},
  numbers=left, numberstyle=\tiny\color{gray}, numbersep=6pt,
  frame=single, breaklines=true, showstringspaces=false,
  columns=fullflexible, keepspaces=true, tabsize=4
}

The three algorithms of Duan et al.~\cite{duan2025} are implemented in Python as methods of
the class \texttt{DuanSSSP}, in the file \texttt{duan.py}: \texttt{find\_pivots}
(Algorithm~1), \texttt{base\_case} (Algorithm~2) and \texttt{bmssp} (Algorithm~3). The
listings below are the code of these methods. Only the lines that write logs and update the
operation counters of the experiment were removed, since they do not change the result. We
checked that this version returns the same distances and predecessors as the full one.

\paragraph{Shared state.} The methods work on four lists indexed by vertex:
\texttt{d[v]} holds the estimate $\hat{d}[v]$, \texttt{h[v]} the number of arcs of the
current path to $v$, \texttt{pred[v]} the predecessor $\mathrm{Pred}[v]$, and
\texttt{pred\_w[v]} the weight of the arc $(\mathrm{Pred}[v], v)$. Adjacency lists are flat
tuples $(v_0, w_0, v_1, w_1, \dots)$, read in pairs with \texttt{zip(adj\_u, adj\_u)}. The
parameter \texttt{depth} is only used to indent the logs.

\paragraph{Total order (Assumption 2.1).} The paper assumes that all paths have distinct
lengths. We represent the value of $\hat{d}[v]$ by the tuple $(\hat{d}[v], h[v], v)$, which
Python compares lexicographically, so every bound $B$, $B_i$ and $B'_i$ is a tuple of this
form ($B=\infty$ is $(\infty,\infty,\infty)$). The relaxation test of Remark~3.4,
$\hat{d}[u]+w_{uv}\le\hat{d}[v]$, becomes
\begin{equation*}
  \bigl(\hat{d}[u]+w,\; h[u]+1,\; u\bigr) \;\le\; \bigl(\hat{d}[v],\; h[v],\; \mathrm{Pred}[v]\bigr),
\end{equation*}
because the two paths compared end at $v$ and the first possible difference after the
length and the number of arcs is the predecessor. The ``$\le$'' lets an arc that is already
tight be relaxed again, as the paper requires. This test is written directly inside each of
the three methods.

\subsubsection{Algorithm 1: FindPivots}

\texttt{find\_pivots(B, S)} (Listing~\ref{lst:findpivots}) decides which vertices of $S$
need to be passed on to the recursion.
\begin{enumerate}
  \item It runs $k$ rounds of Bellman--Ford starting from $S$. In each round only the
        vertices reached in the previous round (\texttt{prev}) are expanded, and a vertex
        enters $W_i$ only if its new value is below $B$.
  \item If $|W|$ exceeds $k|S|$, it stops and returns $P=S$: the frontier is too large to
        be reduced.
  \item Otherwise it builds the forest $F$ of tight arcs inside $W$. Under the total order
        a vertex has at most one tight incoming arc, the one from $\mathrm{Pred}[v]$, so
        $F$ is read from \texttt{pred} in $O(|W|)$ time. The two equalities check that the
        arc is still tight, because $\mathrm{Pred}[v]$ may be outdated if $u$ was improved
        later.
  \item The pivots are the roots in $S$ whose tree in $F$ has at least $k$ vertices. Tree
        sizes are computed with an explicit-stack DFS.
\end{enumerate}

\begin{lstlisting}[style=pycode, caption={Algorithm 1, FindPivots (\texttt{duan.py}).}, label={lst:findpivots}]
def find_pivots(self, B, S, depth):
    k, adj = self.k, self.g.adj
    d, h, pred, pred_w = self.d, self.h, self.pred, self.pred_w
    W = set(S)
    prev = set(S)
    for i in range(1, k + 1):
        Wi = set()
        for u in prev:
            adj_u = iter(adj[u])
            for v, w in zip(adj_u, adj_u):
                du, hu = d[u] + w, h[u] + 1
                if (du, hu, u) <= (d[v], h[v], pred[v]):
                    d[v], h[v], pred[v], pred_w[v] = du, hu, u, w
                    if (du, hu, v) < B:
                        Wi.add(v)
        W |= Wi
        if len(W) > k * len(S):
            return set(S), W
        prev = Wi

    children = {u: [] for u in W}
    has_parent = set()
    for v in W:
        u = pred[v]
        if u in children and d[v] == d[u] + pred_w[v] and h[v] == h[u] + 1:
            children[u].append(v)
            has_parent.add(v)

    P = set()
    for r in S:
        if r in has_parent:
            continue
        size, stack = 0, [r]
        while stack:
            x = stack.pop()
            size += 1
            stack.extend(children[x])
        if size >= k:
            P.add(r)
    return P, W
\end{lstlisting}

\subsubsection{Algorithm 2: BaseCase}

\texttt{base\_case(B, S)} (Listing~\ref{lst:basecase}) is the bottom of the recursion
($l=0$), where $S=\{x\}$ is a single vertex.
\begin{enumerate}
  \item It runs a Dijkstra from $x$ that settles at most $k+1$ vertices and relaxes only
        paths whose value is below $B$.
  \item The paper uses a heap with DecreaseKey. We use Python's binary heap
        (\texttt{heapq}) with lazy deletion: an improved vertex is pushed again, and
        outdated entries are skipped when popped, by comparing the stored key with the
        current one. \texttt{heap\_ops} returns the push and pop functions of the heap
        selected for the run.
  \item If at most $k$ vertices were settled, the execution is successful and returns
        $B'=B$ and all of them. Otherwise the largest value among the $k+1$ settled
        vertices becomes the new bound $B'$, and only the vertices below it are returned.
\end{enumerate}

\begin{lstlisting}[style=pycode, caption={Algorithm 2, BaseCase (\texttt{duan.py}).}, label={lst:basecase}]
def base_case(self, B, S, depth):
    d, h, pred, pred_w = self.d, self.h, self.pred, self.pred_w
    adj = self.g.adj
    (x,) = S
    U0 = {x}
    push, pop = heap_ops(self.contar)
    heap = []
    kx = (d[x], h[x], x)
    push(heap, (kx, x))
    in_heap = {x: kx}
    done = set()
    while heap and len(U0) < self.k + 1:
        kv, u = pop(heap)
        if u in done or kv != (d[u], h[u], u):
            continue
        done.add(u)
        U0.add(u)
        adj_u = iter(adj[u])
        for v, w in zip(adj_u, adj_u):
            du, hu = d[u] + w, h[u] + 1
            if (du, hu, u) <= (d[v], h[v], pred[v]) and (du, hu, v) < B:
                d[v], h[v], pred[v], pred_w[v] = du, hu, u, w
                kv2 = (du, hu, v)
                old = in_heap.get(v)
                if old is None or kv2 < old:
                    in_heap[v] = kv2
                    push(heap, (kv2, v))
    if len(U0) <= self.k:
        return B, U0
    Bp = max((d[v], h[v], v) for v in U0)
    U = {v for v in U0 if (d[v], h[v], v) < Bp}
    return Bp, U
\end{lstlisting}

\subsubsection{Algorithm 3: BMSSP}

\texttt{bmssp(l, B, S)} (Listing~\ref{lst:bmssp}) is the recursive core. The solver starts
with \texttt{bmssp(l, (inf, inf, inf), \{s\})} at the top level $l$.
\begin{enumerate}
  \item At $l=0$ it delegates to \texttt{base\_case}. Otherwise it calls
        \texttt{find\_pivots} and inserts the pivots into the block-based structure $D$ of
        Lemma~3.3 (class \texttt{BlockDS}), with block size $M=2^{(l-1)t}$. $B'_0$ is the
        smallest value among the pivots.
  \item While $|U| < k\cdot 2^{lt}$ and $D$ is not empty, it pulls from $D$ the (up to)
        $M$ smallest keys $S_i$ with a separating bound $B_i$, and recurses with
        \texttt{bmssp(l-1, B\_i, S\_i)}, which returns $(B'_i, U_i)$.
  \item It relaxes the arcs leaving $U_i$. A relaxed vertex whose new value lies in
        $[B_i, B)$ goes back into $D$ with Insert. One in $[B'_i, B_i)$ goes into the list
        \texttt{K}, together with the vertices of $S_i$ that the recursive call did not
        complete. \texttt{K} is returned to $D$ with a single BatchPrepend, since all its
        values are smaller than every value still in $D$.
  \item At the end, $B'=\min(B'_i, B)$, and $U$ also receives the vertices of $W$ (from
        FindPivots) whose value is below $B'$. $B'=B$ means a successful execution;
        $B'<B$ means it stopped early because of a large workload.
\end{enumerate}
The paper checks the workload condition inside the loop; here it is part of the
\texttt{while} condition, which has the same effect.

\begin{lstlisting}[style=pycode, caption={Algorithm 3, BMSSP (\texttt{duan.py}).}, label={lst:bmssp}]
def bmssp(self, l, B, S, depth=0):
    d, h, pred, pred_w = self.d, self.h, self.pred, self.pred_w
    adj = self.g.adj
    if l == 0:
        return self.base_case(B, S, depth)

    P, W = self.find_pivots(B, S, depth)
    M = 2 ** ((l - 1) * self.t)
    D = BlockDS(M, B, self.fmt, depth, self.contar)
    for x in P:
        D.insert(x, (d[x], h[x], x))
    Bp_last = min(((d[x], h[x], x) for x in P), default=B)
    U = set()
    limit = self.k * 2 ** (l * self.t)

    while len(U) < limit and len(D) > 0:
        Bi, Si = D.pull()
        Bpi, Ui = self.bmssp(l - 1, Bi, set(Si), depth + 1)
        U |= Ui

        K = []
        for u in Ui:
            adj_u = iter(adj[u])
            for v, w in zip(adj_u, adj_u):
                du, hu = d[u] + w, h[u] + 1
                if (du, hu, u) <= (d[v], h[v], pred[v]):
                    d[v], h[v], pred[v], pred_w[v] = du, hu, u, w
                    val = (du, hu, v)
                    if Bi <= val < B:
                        D.insert(v, val)
                    elif Bpi <= val < Bi:
                        K.append((v, val))
        K.extend((x, (d[x], h[x], x)) for x in Si if Bpi <= (d[x], h[x], x) < Bi)
        D.batch_prepend(K)
        Bp_last = Bpi

    Bp = min(Bp_last, B)
    U |= {x for x in W if (d[x], h[x], x) < Bp}
    return Bp, U
\end{lstlisting}
 right below that heading.
% Preamble packages required:
%   \usepackage{amsmath}
%   \usepackage{listings}
%   \usepackage{xcolor}
% Citation key used: duan2025 (entry in implementation_details.bib).
% ---------------------------------------------------------------------------

\lstdefinestyle{pycode}{
  language=Python,
  basicstyle=\ttfamily\footnotesize,
  keywordstyle=\bfseries\color{blue!60!black},
  commentstyle=\itshape\color{green!40!black},
  stringstyle=\color{red!60!black},
  numbers=left, numberstyle=\tiny\color{gray}, numbersep=6pt,
  frame=single, breaklines=true, showstringspaces=false,
  columns=fullflexible, keepspaces=true, tabsize=4
}

The three algorithms of Duan et al.~\cite{duan2025} are implemented in Python as methods of
the class \texttt{DuanSSSP}, in the file \texttt{duan.py}: \texttt{find\_pivots}
(Algorithm~1), \texttt{base\_case} (Algorithm~2) and \texttt{bmssp} (Algorithm~3). The
listings below are the code of these methods. Only the lines that write logs and update the
operation counters of the experiment were removed, since they do not change the result. We
checked that this version returns the same distances and predecessors as the full one.

\paragraph{Shared state.} The methods work on four lists indexed by vertex:
\texttt{d[v]} holds the estimate $\hat{d}[v]$, \texttt{h[v]} the number of arcs of the
current path to $v$, \texttt{pred[v]} the predecessor $\mathrm{Pred}[v]$, and
\texttt{pred\_w[v]} the weight of the arc $(\mathrm{Pred}[v], v)$. Adjacency lists are flat
tuples $(v_0, w_0, v_1, w_1, \dots)$, read in pairs with \texttt{zip(adj\_u, adj\_u)}. The
parameter \texttt{depth} is only used to indent the logs.

\paragraph{Total order (Assumption 2.1).} The paper assumes that all paths have distinct
lengths. We represent the value of $\hat{d}[v]$ by the tuple $(\hat{d}[v], h[v], v)$, which
Python compares lexicographically, so every bound $B$, $B_i$ and $B'_i$ is a tuple of this
form ($B=\infty$ is $(\infty,\infty,\infty)$). The relaxation test of Remark~3.4,
$\hat{d}[u]+w_{uv}\le\hat{d}[v]$, becomes
\begin{equation*}
  \bigl(\hat{d}[u]+w,\; h[u]+1,\; u\bigr) \;\le\; \bigl(\hat{d}[v],\; h[v],\; \mathrm{Pred}[v]\bigr),
\end{equation*}
because the two paths compared end at $v$ and the first possible difference after the
length and the number of arcs is the predecessor. The ``$\le$'' lets an arc that is already
tight be relaxed again, as the paper requires. This test is written directly inside each of
the three methods.

\subsubsection{Algorithm 1: FindPivots}

\texttt{find\_pivots(B, S)} (Listing~\ref{lst:findpivots}) decides which vertices of $S$
need to be passed on to the recursion.
\begin{enumerate}
  \item It runs $k$ rounds of Bellman--Ford starting from $S$. In each round only the
        vertices reached in the previous round (\texttt{prev}) are expanded, and a vertex
        enters $W_i$ only if its new value is below $B$.
  \item If $|W|$ exceeds $k|S|$, it stops and returns $P=S$: the frontier is too large to
        be reduced.
  \item Otherwise it builds the forest $F$ of tight arcs inside $W$. Under the total order
        a vertex has at most one tight incoming arc, the one from $\mathrm{Pred}[v]$, so
        $F$ is read from \texttt{pred} in $O(|W|)$ time. The two equalities check that the
        arc is still tight, because $\mathrm{Pred}[v]$ may be outdated if $u$ was improved
        later.
  \item The pivots are the roots in $S$ whose tree in $F$ has at least $k$ vertices. Tree
        sizes are computed with an explicit-stack DFS.
\end{enumerate}

\begin{lstlisting}[style=pycode, caption={Algorithm 1, FindPivots (\texttt{duan.py}).}, label={lst:findpivots}]
def find_pivots(self, B, S, depth):
    k, adj = self.k, self.g.adj
    d, h, pred, pred_w = self.d, self.h, self.pred, self.pred_w
    W = set(S)
    prev = set(S)
    for i in range(1, k + 1):
        Wi = set()
        for u in prev:
            adj_u = iter(adj[u])
            for v, w in zip(adj_u, adj_u):
                du, hu = d[u] + w, h[u] + 1
                if (du, hu, u) <= (d[v], h[v], pred[v]):
                    d[v], h[v], pred[v], pred_w[v] = du, hu, u, w
                    if (du, hu, v) < B:
                        Wi.add(v)
        W |= Wi
        if len(W) > k * len(S):
            return set(S), W
        prev = Wi

    children = {u: [] for u in W}
    has_parent = set()
    for v in W:
        u = pred[v]
        if u in children and d[v] == d[u] + pred_w[v] and h[v] == h[u] + 1:
            children[u].append(v)
            has_parent.add(v)

    P = set()
    for r in S:
        if r in has_parent:
            continue
        size, stack = 0, [r]
        while stack:
            x = stack.pop()
            size += 1
            stack.extend(children[x])
        if size >= k:
            P.add(r)
    return P, W
\end{lstlisting}

\subsubsection{Algorithm 2: BaseCase}

\texttt{base\_case(B, S)} (Listing~\ref{lst:basecase}) is the bottom of the recursion
($l=0$), where $S=\{x\}$ is a single vertex.
\begin{enumerate}
  \item It runs a Dijkstra from $x$ that settles at most $k+1$ vertices and relaxes only
        paths whose value is below $B$.
  \item The paper uses a heap with DecreaseKey. We use Python's binary heap
        (\texttt{heapq}) with lazy deletion: an improved vertex is pushed again, and
        outdated entries are skipped when popped, by comparing the stored key with the
        current one. \texttt{heap\_ops} returns the push and pop functions of the heap
        selected for the run.
  \item If at most $k$ vertices were settled, the execution is successful and returns
        $B'=B$ and all of them. Otherwise the largest value among the $k+1$ settled
        vertices becomes the new bound $B'$, and only the vertices below it are returned.
\end{enumerate}

\begin{lstlisting}[style=pycode, caption={Algorithm 2, BaseCase (\texttt{duan.py}).}, label={lst:basecase}]
def base_case(self, B, S, depth):
    d, h, pred, pred_w = self.d, self.h, self.pred, self.pred_w
    adj = self.g.adj
    (x,) = S
    U0 = {x}
    push, pop = heap_ops(self.contar)
    heap = []
    kx = (d[x], h[x], x)
    push(heap, (kx, x))
    in_heap = {x: kx}
    done = set()
    while heap and len(U0) < self.k + 1:
        kv, u = pop(heap)
        if u in done or kv != (d[u], h[u], u):
            continue
        done.add(u)
        U0.add(u)
        adj_u = iter(adj[u])
        for v, w in zip(adj_u, adj_u):
            du, hu = d[u] + w, h[u] + 1
            if (du, hu, u) <= (d[v], h[v], pred[v]) and (du, hu, v) < B:
                d[v], h[v], pred[v], pred_w[v] = du, hu, u, w
                kv2 = (du, hu, v)
                old = in_heap.get(v)
                if old is None or kv2 < old:
                    in_heap[v] = kv2
                    push(heap, (kv2, v))
    if len(U0) <= self.k:
        return B, U0
    Bp = max((d[v], h[v], v) for v in U0)
    U = {v for v in U0 if (d[v], h[v], v) < Bp}
    return Bp, U
\end{lstlisting}

\subsubsection{Algorithm 3: BMSSP}

\texttt{bmssp(l, B, S)} (Listing~\ref{lst:bmssp}) is the recursive core. The solver starts
with \texttt{bmssp(l, (inf, inf, inf), \{s\})} at the top level $l$.
\begin{enumerate}
  \item At $l=0$ it delegates to \texttt{base\_case}. Otherwise it calls
        \texttt{find\_pivots} and inserts the pivots into the block-based structure $D$ of
        Lemma~3.3 (class \texttt{BlockDS}), with block size $M=2^{(l-1)t}$. $B'_0$ is the
        smallest value among the pivots.
  \item While $|U| < k\cdot 2^{lt}$ and $D$ is not empty, it pulls from $D$ the (up to)
        $M$ smallest keys $S_i$ with a separating bound $B_i$, and recurses with
        \texttt{bmssp(l-1, B\_i, S\_i)}, which returns $(B'_i, U_i)$.
  \item It relaxes the arcs leaving $U_i$. A relaxed vertex whose new value lies in
        $[B_i, B)$ goes back into $D$ with Insert. One in $[B'_i, B_i)$ goes into the list
        \texttt{K}, together with the vertices of $S_i$ that the recursive call did not
        complete. \texttt{K} is returned to $D$ with a single BatchPrepend, since all its
        values are smaller than every value still in $D$.
  \item At the end, $B'=\min(B'_i, B)$, and $U$ also receives the vertices of $W$ (from
        FindPivots) whose value is below $B'$. $B'=B$ means a successful execution;
        $B'<B$ means it stopped early because of a large workload.
\end{enumerate}
The paper checks the workload condition inside the loop; here it is part of the
\texttt{while} condition, which has the same effect.

\begin{lstlisting}[style=pycode, caption={Algorithm 3, BMSSP (\texttt{duan.py}).}, label={lst:bmssp}]
def bmssp(self, l, B, S, depth=0):
    d, h, pred, pred_w = self.d, self.h, self.pred, self.pred_w
    adj = self.g.adj
    if l == 0:
        return self.base_case(B, S, depth)

    P, W = self.find_pivots(B, S, depth)
    M = 2 ** ((l - 1) * self.t)
    D = BlockDS(M, B, self.fmt, depth, self.contar)
    for x in P:
        D.insert(x, (d[x], h[x], x))
    Bp_last = min(((d[x], h[x], x) for x in P), default=B)
    U = set()
    limit = self.k * 2 ** (l * self.t)

    while len(U) < limit and len(D) > 0:
        Bi, Si = D.pull()
        Bpi, Ui = self.bmssp(l - 1, Bi, set(Si), depth + 1)
        U |= Ui

        K = []
        for u in Ui:
            adj_u = iter(adj[u])
            for v, w in zip(adj_u, adj_u):
                du, hu = d[u] + w, h[u] + 1
                if (du, hu, u) <= (d[v], h[v], pred[v]):
                    d[v], h[v], pred[v], pred_w[v] = du, hu, u, w
                    val = (du, hu, v)
                    if Bi <= val < B:
                        D.insert(v, val)
                    elif Bpi <= val < Bi:
                        K.append((v, val))
        K.extend((x, (d[x], h[x], x)) for x in Si if Bpi <= (d[x], h[x], x) < Bi)
        D.batch_prepend(K)
        Bp_last = Bpi

    Bp = min(Bp_last, B)
    U |= {x for x in W if (d[x], h[x], x) < Bp}
    return Bp, U
\end{lstlisting}
